\subsection{A certified cyclotomic base change and its six local norm characters}

\begin{statement}{E.76}{Proved}{The first two genus quotients after adjoining a ninth root of unity}
Keep $K$, $L_d$, $\beta_d$, and $\mathfrak q_{3,d}$ from E.72.
Choose $\xi^3=\zeta$, with $\xi$ of order nine, and put
\[
 B_d=L_d(\xi),\qquad B=B_1
   =\mathbb Q(\xi,a),\qquad a^3=17,\quad\beta_1=a\xi.
 \tag{E76.1}
\]
For every $d\ge1$, $B_d/L_d$ is cyclic of degree three,
unramified at every finite prime, and split completely at the
unique primes above $3$ and $17$. Thus $B_d$ has three primes
$P_{d,j}$ above $3$ and three primes $Q_{d,j}$ above $17$.
The adjacent extensions $B_{d+1}/B_d$ are totally ramified
at exactly these six primes. The extensions $B_{d+2}/B_d$
are cyclic of degree nine and are totally ramified at the same
six primes.

The full ideal class group of $B$, with the certified coordinates
specified below, is
\[
 \operatorname{Cl}(B)=(\mathbb Z/3)^2,\qquad
 [P_{1,j}]=(1,0),\quad[Q_{1,j}]=(2,0)\quad(1\le j\le3).
 \tag{E76.2}
\]
The second class generator is the explicit prime of norm $109$
in the certificate. In particular the six ramified prime classes
generate precisely the first cyclic summand.

Let $U=\mathcal O_B^\times$, and order the six local places as
$P_{1,1},P_{1,2},P_{1,3},Q_{1,1},Q_{1,2},Q_{1,3}$.
Let $(\xi,u_1,\ldots,u_8)$ be the complete certified unit
generators in \texttt{e\_basechange\_field.gp}; the omitted
generator $-1$ is a norm in both odd-degree extensions.
The six local unit-norm quotient characters for $B_2/B$,
in the explicitly recorded local generators, have matrix
\[
 M_3=\begin{pmatrix}
 2&1&0&1&0&0&0&2&1\\
 2&1&0&1&0&0&0&0&0\\
 2&1&0&1&0&0&0&1&2\\
 2&1&0&1&0&0&0&1&1\\
 2&1&0&1&0&0&0&1&1\\
 1&2&0&2&0&0&0&2&2
 \end{pmatrix}\pmod3.
 \tag{E76.3}
\]
Its rank is three. For $B_3/B$, the corresponding six
characters take values in $\mathbb Z/9$ and have matrix
\[
 M_9=\begin{pmatrix}
 1&5&3&2&6&0&3&4&2\\
 1&5&3&8&3&6&0&3&6\\
 1&5&3&5&0&3&6&8&7\\
 5&7&6&4&0&0&3&1&4\\
 8&4&6&7&3&3&0&1&1\\
 1&5&3&5&3&3&3&5&8
 \end{pmatrix}\pmod9.
 \tag{E76.4}
\]
Here
$\operatorname{coker}(M_9:(\mathbb Z/9)^9\to(\mathbb Z/9)^6)
 \simeq(\mathbb Z/9)^2\oplus\mathbb Z/3$.
The displayed maps are actual local norm-quotient coordinates;
they do not identify the six quotient generators with a single
common Hilbert-symbol root. Their kernels, the simultaneous
kernel, and the following indices are consequently fixed
without a Hilbert-symbol sign convention:
\[
 [U:U\cap N_{B_2/B}B_2^\times]=27,\qquad
 [U:U\cap N_{B_3/B}B_3^\times]=2187.
 \tag{E76.5}
\]
For the actual cyclic groups $G_n=\operatorname{Gal}(B_{1+\log_3n}/B)$,
$n=3,9$, these give
\[
 |\operatorname{Cl}(B_2)^{G_3}|=81,\qquad
 |\operatorname{Cl}(B_3)^{G_9}|=243.
 \tag{E76.6}
\]

There are more precise finite quotient maps. Let $\mathcal G_n/B$
be the largest abelian extension containing $B_{1+\log_3n}$
that is unramified over that field. This is a finite subfield of
its Hilbert class field. Write $\Gamma_n=\operatorname{Gal}(\mathcal G_n/B)$.
The ray-class presentations and their full integral Smith
coordinate maps give
\[
 \Gamma_3\simeq\mathbb Z/9\oplus(\mathbb Z/3)^3,
 \qquad
 \Gamma_9\simeq(\mathbb Z/9)^3\oplus\mathbb Z/3.
 \tag{E76.7}
\]
For each $j$, the Frobenius of $P_{2,j}$ in
$\mathcal G_3/B_2$ has Smith coordinates $(3,1,0,1)$.
The Frobenii of $P_{3,1},P_{3,2},P_{3,3}$ in
$\mathcal G_9/B_3$ have respective coordinates
\[
                  (3,6,0,2),\quad(6,6,6,2),\quad(0,6,3,2).
 \tag{E76.8}
\]
Every one of these six elements has exact order three.
The inclusion $\mathcal G_3\subset\mathcal G_9$ induces
the explicit surjective matrix in the certificate; it sends
each prime Frobenius in (E76.8) to its counterpart in
(E76.7).

The actual automorphism
\[
 \delta(\xi)=\xi^4,\qquad\delta(a)=a\xi^{-3},
 \qquad\delta(\beta_1)=\beta_1
 \tag{E76.9}
\]
acts on $\Gamma_9$. The coordinate calculation gives
$\Gamma_9/(\delta-1)\Gamma_9\simeq(\mathbb Z/3)^3$.
Let $A_9=\operatorname{Gal}(\mathcal G_9/B_3)$, the kernel
of the actual Kummer character $\Gamma_9\to\mathbb Z/9$.
Then
\[
 A_9\simeq(\mathbb Z/9)^2\oplus\mathbb Z/3,
 \qquad A_9/(\delta-1)A_9\simeq(\mathbb Z/3)^2.
 \tag{E76.10}
\]
In the latter Smith coordinates all three elements of
(E76.8) have image $(2,0)$.

The comparison with E.72 uses the literal ideal maps
\[
 \begin{split}
 N_{B_d/L_d}[P_{d,j}]&=[\mathfrak q_{3,d}],\\
 j_d[\mathfrak q_{3,d}]&=\sum_{j=1}^3[P_{d,j}],\\
 N_{B_d/L_d}j_d&=3,\qquad
 j_dN_{B_d/L_d}=1+\delta+\delta^2 .
 \end{split}
 \tag{E76.11}
\]
Moreover $\mathcal G_9/L_3$ is unramified and Galois.
The subgroup $A_9/(\delta-1)A_9$ injects into
$\operatorname{Gal}(\mathcal G_9/L_3)^{\mathrm{ab}}$;
the class $[\mathfrak q_{3,3}]$ maps to $(2,0)$ in this
subgroup. This gives the explicit order-three quotient of
that original prime class. E.72 still gives the full bound
$3\mid\operatorname{ord}[\mathfrak q_{3,3}]\mid9$.
The finite quotients here evaluate the stated maps, without
evaluating the remaining order in $\operatorname{Cl}(L_3)$
or the eventual integer $h$.

Finally $|\operatorname{Cl}(B_2)\{3\}|\ge81>9
=|\operatorname{Cl}(B)\{3\}|$. Thus the initial norm map
on three-class groups is surjective with nontrivial kernel.
This is an exact failed initial equality for the stability
test of Li--Ouyang--Xu--Zhang, Proposition~3.2, applied to
$B_3/B_2/B$; all its cyclicity and ramification requirements
hold. The computation does not supply a later equality.
The source motivation remains LM~\S51.16,
fols.~733v--736v, scans~172--178: these are concrete
arithmetic calculations inside the norm systems of
E.46 and E.72.
\end{statement}

\begin{proof}
First fix all the field maps. The element $17$ has a cube
root in $\mathbb Q_3$: the cube map from $1+3\mathbb Z_3$
to $1+9\mathbb Z_3$ is bijective by the three-adic logarithm,
and $-17\in1+9\mathbb Z_3$. Hence $\xi=\beta_1/a$ belongs
to the completion of $L_1$ at three. At seventeen, the
completion of $K$ is the unramified quadratic extension
of $\mathbb Q_{17}$ and its residue field has order $289$.
It contains $\mu_9$ by Teichmuller lifting. The extension
$K(\xi)/K$ is unramified away from three. These local maps
show that its base change to every $L_d$ is everywhere
unramified and split at three and seventeen. The degree
remains three: $L_d/K$ is totally ramified at its unique
prime above seventeen, whereas seventeen splits into
three primes in $K(\xi)/K$, excluding an inclusion
$K(\xi)\subset L_d$.

The element $\beta_d+1$ is a uniformizer at three.
Indeed the polynomial $(-1+Y)^{3^d}-17\zeta$ is
Eisenstein over $K_3$: its constant coefficient
$-1-17\zeta$ has valuation one and its intermediate
coefficients are divisible by three. The corresponding
assertion at seventeen follows directly from
$Y^{3^d}-17\zeta$. The split completions above prove
the stated total ramification in the $B_d$ chain.
The polynomial $Y^9-\beta_d$ defines $B_{d+2}/B_d$,
and $\mu_9\subset B_d$ makes this extension cyclic.

Here is the unconditional class and unit certificate for
$B$. The chosen primitive element is $t=a-\xi$, with
minimal polynomial
\[
 t^{18}-105t^{15}+4137t^{12}-81182t^9
              +1128687t^6-6822222t^3+28934443.
 \tag{E76.12}
\]
The exact maps from $\xi,a$ to this polynomial field,
the integral basis, both class-generator ideals and all
fundamental units are retained. PARI~2.19.0 computes
the discriminant $-3^{27}17^{12}$ and returns
\texttt{bnfcertify}=1 for the full class group and unit
basis. That return value certifies the calculation
without a generalized Riemann hypothesis. The six
\texttt{bnfisprincipal} coordinates give (E76.2);
the first generator is $P_{1,1}$ and the other is the
retained prime above $109$. These are exact ideals in
the displayed integral basis, not classes inferred
from an unsuccessful principal-element search.
The primary software descriptions and the complete
fresh execution are part of the certificate
\cite{E_PARICertifiedUnits}.

We explain why the finite local computations determine
the entire six characters. At a place $P_{1,j}$ the
field $B_{P_{1,j}}$ is $\mathbb Q_3(\xi)$, of absolute
ramification degree six. Write $U_k=1+P_{1,j}^k$.
Logarithm and exponential are inverse on $U_4$, so
\[
                     U_{10}=U_4^3,\qquad U_{16}=U_4^9.
 \tag{E76.13}
\]
Thus reduction modulo $P^{10}$, respectively $P^{16}$,
followed by quotienting by third, respectively ninth,
powers computes the actual local unit quotient by
powers. Its elementary divisors are seven copies of
$3$, respectively $9$. At the tame places the
corresponding quotients are the third- and ninth-power
quotients of $\mathbb F_{289}^\times$, each cyclic of
the indicated order.

For $n=3,9$ put $Y^n=\beta_1$. The element $Y+1$
is a uniformizer of the local degree-$n$ extension.
For each $k$ the program computes the literal
multiplication determinant
\[
                  N_k=N(1+(Y+1)^k)
 \tag{E76.14}
\]
in $B[Y]/(Y^n-\beta_1)$. All these elements are
local units at the three wild places. For $n=3$,
the columns with $k=1,2,4,5,7,10$ have rank six
modulo three. Their annihilator in the recorded
seven local coordinates is
$(2,0,2,1,2,1,0)$. For $n=9$, the first ten columns
together with $9I_7$ have Smith quotient $\mathbb Z/9$
at each wild place. The respective quotient rows are
\[
 \begin{gathered}
 \begin{pmatrix}1&-8&-4&-7&-1&-6&-3\end{pmatrix},\qquad
 \begin{pmatrix}1&-8&-7&-1&-4&-6&-3\end{pmatrix},\\[3pt]
 \begin{pmatrix}1&-8&-1&-4&-7&-6&-3\end{pmatrix}.
 \end{gathered}
 \tag{E76.15}
\]
Local reciprocity identifies the unit norm quotient
of a totally ramified cyclic degree-$n$ extension
with $\mathbb Z/n$. Consequently each computed
subgroup of norm units, already of index $n$, is
the entire unit norm kernel. At seventeen the tame
unit norm kernel is exactly the subgroup with
$n$th-power residue. This proves the characters,
including their precise coordinate meanings
\cite[III, Theorems~3.1, 3.4]{E_MilneCFT}.

Applying these six quotient rows to the complete
global generators gives (E76.3)--(E76.4). The
integer Smith calculations include their unimodular
transformations and inverse maps. A global unit
is automatically a local norm at every other
place, since the extension there is unramified.
The Hasse norm theorem identifies the simultaneous
kernel with $U\cap N B_{1+\log_3n}^{\times}$
\cite[VIII, Theorem~3.1]{E_MilneCFT}.
Its index is $3^3$ for $n=3$ and
$9^6/(9^2\cdot3)=3^7$ for $n=9$.
Chevalley's exact ambiguous class formula now reads
\[
 |\operatorname{Cl}(B_{1+\log_3n})^{G_n}|
   ={9\,n^6\over n\,[U:U\cap N B_{1+\log_3n}^{\times}]},
 \tag{E76.16}
\]
which gives (E76.6). The degree-three formula and
its complete ideal/unit proof are
\cite[Theorem~1]{E_LemmermeyerAmbiguous}; the general
cyclic formula is also stated explicitly in
\cite[Theorem~2.6, Remark~2.7(2), Lemma~2.8]{E_LOXZKummer}.

We next identify the finite ray quotients with the
claimed fields, including the relevant quotient
maps. Let $\mathfrak m_n$ have exponent ten at the
three wild primes for $n=3$, exponent sixteen for
$n=9$, and exponent one at the three tame primes.
In $\operatorname{Cl}_{\mathfrak m_n}(B)$ quotient
by all local norm-unit subgroups. The local maps
are realized by the Chinese remainder theorem:
put a chosen local norm unit at one place and
put $1$ at the other five, and take the ray class
of its principal ideal. These are exactly the
columns in the retained presentation, together
with the ray cyclic relations. The finite idelic
description is
\[
 \Gamma_n=\mathbb A_{B,f}^{\times}\big/
 B^\times\left(
 \prod_{v\in S}N\mathcal O_{B_{1+\log_3n},w}^{\times}
 \prod_{v\notin S}\mathcal O_{B,v}^{\times}\right).
 \tag{E76.17}
\]
The global norm subgroup for the cyclic extension
contains the denominator, so its class field
contains $B_{1+\log_3n}$. At each ramified place,
the inertia quotient has size at most $n$ and
maps onto the inertia group of that cyclic field,
of size $n$. It follows that the resulting extension
over that field is unramified. Conversely local
norm functoriality kills precisely the displayed
norm units in every abelian extension of $B$
unramified over $B_{1+\log_3n}$. This proves the
maximal property defining $\mathcal G_n$ by the
actual Artin map
\cite[III, Lemma~3.2; V, Proposition~5.2]{E_MilneCFT}.

The norm compatibility used here follows from the proved
cup-product comparison in \cite[III, Theorem~3.1 and its
restriction/corestriction proof]{E_MilneCFT}.
For each actual finite local tower $L/E/K$ above, put
$G=\operatorname{Gal}(L/K)$ and $H=\operatorname{Gal}(L/E)$.
Corestriction in Tate degree $-2$ is the inclusion-induced map
$H^{\mathrm{ab}}\to G^{\mathrm{ab}}$, and in degree zero
it is $N_{E/K}$ on the multiplicative norm quotients.
The fundamental classes satisfy
$\operatorname{Res}_H^G u_{L/K}=u_{L/E}$.
The projection formula
\[
 \operatorname{Cor}_H^G(x\smile u_{L/E})
   =\operatorname{Cor}_H^G(x)\smile u_{L/K}
\]
therefore says that the inverse Artin maps commute with
these two maps. Inverting those isomorphisms gives
$\operatorname{Art}_K(N_{E/K}z)
 =\iota\operatorname{Art}_E(z)$, where
$\iota:H^{\mathrm{ab}}\to G^{\mathrm{ab}}$ is induced
by inclusion. This proves precisely the norm and local-unit
compatibilities used in the genus and prime-Frobenius maps.

This also gives a direct proof of the ambiguous-order
formula in the two cases used here. The kernel of the
surjection from (E76.17) to $\operatorname{Cl}(B)$ is
the quotient of the six local unit norm quotients by
the image of $U$. Consequently its order is
$n^6/[U:U\cap N B_{1+\log_3n}^{\times}]$.
Let $H$ be the full Hilbert class field of
$E=B_{1+\log_3n}$. The group
$\operatorname{Gal}(H/B)$ is generated by the abelian
subgroup $\operatorname{Cl}(E)$ and one lift of a
generator $\sigma$ of $\operatorname{Gal}(E/B)$.
Its commutator subgroup is exactly
$(\sigma-1)\operatorname{Cl}(E)$. Its maximal
abelian quotient is $\Gamma_n$, with kernel over
$\operatorname{Gal}(E/B)$ equal to
$\operatorname{Cl}(E)/(\sigma-1)\operatorname{Cl}(E)$.
For a finite abelian group the kernel and cokernel
of $\sigma-1$ have equal order. Dividing
$|\Gamma_n|$ by $n$ therefore proves (E76.16)
and (E76.6) directly from these exact maps.

The quotient in (E76.17) is an extension of the
ordinary class group $(\mathbb Z/3)^2$ by a
quotient of $(\mathbb Z/n)^6$. Its exponent
therefore divides $3n$. This proves that using
the ray class group modulo $3n$th powers in
the certificate loses no element of $\Gamma_n$.
Its relation matrices have exact integer Smith
forms (E76.7). The local norm groups for $n=9$
are contained in those for $n=3$. This proves
$\mathcal G_3\subset\mathcal G_9$ and the
surjection between the displayed ray presentations.
The program computes this map on actual ray ideals;
its change of ray generators is recorded with an
inverse modulo the cyclic relation lattices.

For clarity, the prime Frobenii are computed by
a local uniformizer norm, not by assigning an
unramified symbol to a ramified prime of $B$.
For both values of $n$,
\[
          N(Y+1)=\beta_1+1=\pi.
 \tag{E76.18}
\]
For the chosen wild prime $P$, take $x\equiv\pi$
modulo $P^{1+v_P(\mathfrak m_n)}$ and
$x\equiv1$ at the other modulus factors.
The ideal $I=(x)P^{-1}$ is coprime to the modulus.
Global reciprocity identifies the local Artin
class of $\pi$ with $-[I]$ in additive ray
coordinates. Norm compatibility identifies it
with the Frobenius of the upper prime in
$\mathcal G_n/B_{1+\log_3n}$. Applying the Smith
maps gives exactly (E76.8) and its cubic counterpart.

The final descent uses the actual automorphism,
not an abstract permutation of Smith factors.
In the primitive element (E76.12), it is
$t\mapsto a\xi^{-3}-\xi^4$. The program verifies
the polynomial relation, $\delta(\beta_1)=\beta_1$,
and $\delta(\xi)=\xi^4$. Its action matrix preserves
the ray relation lattice and the local norm-unit
lattice, giving the stated action on $\Gamma_9$.
The Kummer character $\Gamma_9\to\mathbb Z/9$
is evaluated on actual ideal generators: at an
unramified prime $\mathfrak l$, its coordinate is
the unique $e\pmod9$ with
\[
       \overline{\beta_1}^{\,(N\mathfrak l-1)/9}
                         =\overline\xi^{\,e}.
 \tag{E76.19}
\]
In the retained sixteen ray coordinates the row is
\[
 (8,7,5,8,4,5,2,8,3,2,8,5,3,8,1,0)\pmod9.
 \tag{E76.20}
\]
The exact lattice kernel of this row, the
restriction of $\delta$ to it, and both Smith
coordinate maps prove (E76.10). They give image
$(2,0)$ for each of the three prime Frobenii.

The extension $B_3/L_3$ is generated by the
restriction of $\delta$, extending (E76.9) by
$\delta(\beta_3)=\beta_3$. Its conjugation on
$\operatorname{Gal}(B_3/B)$ raises a generator
to the fourth power, so $\mathcal G_9$ is stable
under it by the maximal property already proved.
Hence $\mathcal G_9/L_3$ is Galois and unramified.
For the exact sequence with abelian kernel
\[
 1\longrightarrow A_9\longrightarrow
 \operatorname{Gal}(\mathcal G_9/L_3)
 \longrightarrow\mathbb Z/3\longrightarrow1,
 \tag{E76.21}
\]
the commutator subgroup is exactly $(\delta-1)A_9$:
the group is generated by $A_9$ and one lift of
the cyclic quotient, and its quotient by these
commutators is abelian. This proves the claimed
injection into the abelianization. Since
$\mathfrak q_{3,3}$ splits in $B_3/L_3$, its
Frobenius belongs to this kernel and is the
computed upper Frobenius. This establishes the
precise transfer back to the original prime.

Finally (E76.11) follows directly from residue
degree one and complete splitting, and holds
on full class groups and on their three-primary
parts. The norm from $B_2$ to $B$ is surjective:
its cyclic degree-three extension is ramified,
so it has no nontrivial unramified subextension;
write $H_B$ for the Hilbert class field of $B$.
Then $H_B\cap B_2=B$, and $H_BB_2/B_2$ is unramified
and abelian. Its Artin quotient gives the surjection
\[
 \operatorname{Cl}(B_2)\twoheadrightarrow
 \operatorname{Gal}(H_BB_2/B_2)
 \xrightarrow{\ \mathrm{res}\ }\operatorname{Gal}(H_B/B)
 \simeq\operatorname{Cl}(B).
\]
Norm functoriality identifies this actual composite with
$N_{B_2/B}$. Equation (E76.6) then shows that
the norm kernel has order at least nine. Thus
the equality premise in the opened stability
theorem fails at this actual first step
\cite[Proposition~3.2 and its full proof]{E_LOXZKummer}.
The classes in the kernels of the displayed
genus maps have not been evaluated here; the
statement retains exactly the remaining
prime-class orders from E.72.
\end{proof}
